\section{Calibration}
\label{sec:calibration}

This section reports the headline estimand. It is organised around the two system signals,
because the sweep's central finding is that they behave differently and that conflating
them produces contradictory conclusions.

\subsection{The two deltas}

Figure~\ref{fig:delta-bench-topo} shows both deltas by benchmark and topology, and
Appendix~\ref{app:marginals} gives the full marginal tables.

\begin{figure}[H]
\centering
\includegraphics[width=\linewidth,height=0.32\textheight,keepaspectratio]{fig_delta_by_bench_topo.pdf}
\caption{Coordination calibration penalty by benchmark and topology. Left: the
vote-fraction signal. Right: the final-producer signal. Note the difference in vertical
scale between the two panels.}
\label{fig:delta-bench-topo}
\end{figure}

The vote-fraction delta is large and its sign depends on the task family. On GPQA it runs
from $+0.1508$ for independent to $+0.1897$ for centralized. On TruthfulQA it runs from
$+0.0474$ to $+0.0906$. On MMLU-Pro it is negative throughout, from $-0.1554$ for
centralized to $-0.1860$ for decentralized.

The underlying calibration errors explain the reversal. On GPQA the per-agent error is
0.234 to 0.240 while the vote error is 0.385 to 0.430: agreement is far more confident than
it should be on a benchmark where three agents frequently agree on a wrong answer. On
MMLU-Pro the per-agent error is 0.393 to 0.408 while the vote error is 0.216 to 0.253:
agreement is conservative, because with up to ten options unanimous agreement is rarer and
therefore more informative. The same three-atom signal is overconfident on one task family
and underconfident on another, and neither the signal nor the aggregation rule changed
between them.

The final-producer delta is an order of magnitude smaller and consistently positive. It
runs from $+0.0037$ to $+0.0204$ across all nine benchmark-topology combinations where it
is defined, with the largest value on MMLU-Pro under decentralized debate. Coordination
therefore barely degrades the calibration of the model whose output determines the answer,
even where it badly degrades the honesty of the agreement statistic.

This separation is the report's principal calibration result and it has a direct
implication. A system that displays agreement as confidence can be badly miscalibrated in
either direction depending on the task. A system that displays the decisive model's own
probability incurs a penalty of roughly 0.4 to 2.0 calibration points, which is small
enough to be corrected post hoc. Most of the headline $\dece$ is a statement about how
agreement is converted into a number, not about coordination corrupting the underlying
models.

The single-agent rows in the same table illustrate why the vote signal must not be read
across the agent-count boundary. With one agent the vote fraction is identically one, so
its calibration error is one minus the accuracy by construction: 0.5997 on GPQA, 0.3877 on
MMLU-Pro, 0.3751 on TruthfulQA. These are not model-calibration measurements and the
one-agent vote delta should not be interpreted.

\subsection{Capacity reverses the expected direction}

The sweep was designed to test whether the coordination calibration penalty grows with
capacity. It does not. Fitting power laws to calibration error against parameter count and
contrasting the exponents gives $b_{\mathrm{pa}} - b_{\mathrm{vote}}$ of $+0.077$ $[0.031,
0.127]$ on GPQA, $+0.495$ $[0.424, 0.567]$ on MMLU-Pro and $+0.087$ $[0.008, 0.153]$ on
TruthfulQA. All three intervals exclude zero and all three are positive, meaning vote
calibration error falls faster with scale than per-agent calibration error does.

The consequence is that $\dvote$ shrinks with capacity rather than growing. On the two
largest models the vote delta turns negative on several benchmark and topology
combinations, meaning coordination begins to correct miscalibration rather than amplify it.

Figure~\ref{fig:ece-vs-accuracy} shows the joint distribution of accuracy and calibration
error by capacity for all three signals.

\begin{figure}[H]
\centering
\includegraphics[width=\linewidth,height=0.30\textheight,keepaspectratio]{fig_ece_vs_accuracy.pdf}
\caption{Calibration error against cell accuracy by capacity, for the per-agent, vote and
final-producer signals. The vote panel shows the strongest capacity ordering.}
\label{fig:ece-vs-accuracy}
\end{figure}

Two caveats bound this result. The vote signal has only three atoms with three agents, so
part of the capacity effect could be a change in how mass distributes over a coarse support
rather than a change in the honesty of agreement. Distinguishing these requires the same
measurement at five and seven agents, which the sibling runs are designed to provide. And
the final-producer contrast, which does not suffer from atom coarseness, shows no
comparable capacity effect, which is consistent with the coarseness explanation.

\subsection{Reasoning shifts the confidence distribution rather than sharpening it}

Section~\ref{sec:marginals} reported that reasoning improves accuracy, reduces signed
overconfidence, and simultaneously increases per-agent calibration error. The three
together are only consistent if reasoning moves confidence past the calibration point.

The signed-gap slopes are negative, at $-0.0082$, $-0.0140$ and $-0.0058$ per unit
log-dose on the three multiple-choice benchmarks, so the average confidence falls relative
to accuracy. The per-agent calibration slopes are positive, at $+0.0045$, $+0.0126$ and
$+0.0034$. A model that started overconfident and whose mean confidence is being pushed
down would first improve and then worsen as it passes through calibration into
underconfidence, and the observed combination of a falling gap and a rising error is the
second half of that trajectory.

The design implication is stated in Section~\ref{sec:interpretation}: a system that uses
extended reasoning should not expose raw option logprobs as probabilities without
recalibration, because the direction of the miscalibration has changed even though its
magnitude at the mean has not obviously worsened.

\subsection{The confidence-signal tournament}

Six candidate confidence signals were compared by calibration error within each stratum of
topology, capacity band and reasoning anchor, giving eighteen strata.

The winner is the mean per-agent confidence, which takes 13 of 18 strata. Mean agreeing
verbalized confidence takes the remaining 5. The mean agreeing option-logprob confidence
wins zero of 18.

That last result is the one worth attention, because mean agreeing logprob is the signal a
system would most naturally construct: take the agents that agreed and average their
probabilities. It is beaten in every stratum by simply averaging over all agents without
conditioning on agreement. Conditioning on the winning side selects for the agents that
were confident in that answer, which inflates the average without adding information.

The tournament result carries one important qualification. The winning signal, mean
per-agent confidence, is itself an ensembled quantity and is not the per-agent baseline
used in Equation~\ref{eq:estimand}. It is a system-level signal built from per-agent
inputs. The tournament therefore says the best available system confidence is a simple
average over all agents, not that no aggregation is worthwhile.

Within centralized cells specifically, the orchestrator's own logprob confidence beats the
vote fraction in 4 of 6 strata, which is a weak advantage and not a reliable one.

\subsection{Voting-lift decomposition and error correlation}

If agents made independent errors, the accuracy of a majority vote would be a computable
function of individual accuracy. The gap between the observed vote accuracy and that
independence bound is the correlation penalty.

The decomposition finds no established penalty. The within-cell plug-in penalty has an
interval below zero in 0 of 24 strata, and the observed team vote falls below the
cross-seed independence bound in 1 of 24. On this evidence the sweep cannot demonstrate
that error correlation reduces voting lift.

This is a negative result with a known cause rather than a clean absence of effect. Option
letters are shuffled independently per question and per seed, so the same wrong answer
receives different letters in different cells. Cross-seed wrong-answer collision rates are
therefore not identifiable from the recorded data: two agents choosing the same incorrect
content in different cells appear as different letters. Establishing or refuting the
correlation penalty requires logging canonical answer identities rather than positional
letters, which is recorded as a follow-up in Section~\ref{sec:limitations}.

The item-level evidence in Section~\ref{sec:items} bears on the same mechanism from a
different direction and does find an effect, which is why the two results are not in
conflict.

\subsection{Reliability structure}

The atlas figures in Appendix~\ref{app:atlas} include reliability diagrams for per-agent
and system confidence, faceted by capacity band and reasoning rung, and the Cox
recalibration intercepts and slopes are reported per cell in the marginal tables.

The Cox slope is the relevant summary for the reasoning result: a slope below one indicates
confidence that is too sharp for its discriminative content, and a slope above one
indicates the reverse. Together with the signed gap it distinguishes a signal that is
biased from one that is insufficiently sharp, which the calibration error alone does not.

Post-hoc temperature scaling is also fitted per cell, and provides the ceiling on what
recalibration could recover. The gap between the raw calibration error and the
temperature-scaled error is the portion of miscalibration that is a simple monotone
distortion rather than a failure of ranking.

\subsection{The mathematical benchmark}

The mathematical benchmark has no option-logprob path, so per-agent and final-producer
calibration are undefined and only agreement-based and sampling-based signals survive. The
vote-fraction calibration error is still measurable and is reported throughout.

This makes the benchmark a useful control on the whole calibration analysis. Every result
that depends on option logprobs disappears there, and every result that survives is one
that a deployed free-form system could actually reproduce. The self-consistency and
semantic-entropy signals computed for single-agent cells are the only confidence families
available in that setting, and their behaviour is reported in the atlas.
